% TOP_PROBLEM: {"id": 8700023, "title": "Conjecture 3.5 — Strong Five Exponentials Conjecture", "category_id": 3, "category": {"id": 3, "name": "graph_theory", "display_name": "Graph Theory", "description": "Problems involving graphs, networks, and their properties.", "slug": "graph-theory", "order_index": 3, "created_at": "2026-07-31T15:26:25.671Z"}, "status": "open", "problem_number": "AMR-086-0023", "set_id": 13}
% FIRST_POSTED: 2026-10-01
% ENTRY_KIND: statement_only
% UnsolvedMath ID 8700023; source catalogue code AMR-086-0023
% !TeX program = lualatex
% Definition and sources only; this file is not an LLM solution attempt.
\documentclass[11pt,a4paper]{article}
\usepackage[margin=25mm]{geometry}
\usepackage{fontspec}
\setmainfont{lmroman10-regular.otf}[BoldFont=lmroman10-bold.otf,
 ItalicFont=lmroman10-italic.otf,BoldItalicFont=lmroman10-bolditalic.otf]
\usepackage{amsmath,amssymb,mathtools}
\usepackage{unicode-math}
\setmathfont{latinmodern-math.otf}
\AtBeginDocument{\let\setminus\smallsetminus}
\usepackage{xurl,hyperref}
\hypersetup{colorlinks=true,urlcolor=blue}
\setcounter{secnumdepth}{0}
\setlength{\parindent}{0pt}
\setlength{\parskip}{5pt}
\setlength{\emergencystretch}{3em}
\begin{document}
\section*{Conjecture 3.5 — Strong Five Exponentials Conjecture}\label{8700023}
{\small UnsolvedMath ID 8700023 · AMR-086-0023 · Graph Theory · AMR Open Problem Lists}
\subsection{Definitions and mathematical statement}
Choose $x_1,x_2,y_1,y_2\in\mathbb C$ such that each pair $(x_1,x_2)$ and $(y_1,y_2)$ is linearly independent over $\mathbb Q$. This means that a rational linear relation within either pair has both coefficients zero; in particular, $x_2\ne0$. Let $\beta_{11},\beta_{12},\beta_{21},\beta_{22},\gamma_1,\gamma_2$ be algebraic complex numbers with $\gamma_1\ne0$.

Assume that all five numbers
\[
\exp(x_i y_j-\beta_{ij})\quad(1\le i,j\le2),
\qquad \exp(\gamma_1x_1/x_2-\gamma_2)
\]
are algebraic. Must the five corresponding exponents all be zero? In full, the requested conclusion is
\[
x_i y_j=\beta_{ij}\ (1\le i,j\le2),\qquad
\gamma_1x_1=\gamma_2x_2.
\]
All exponentials here are the ordinary entire complex exponential function; no choice of logarithm or power branch enters their definition.

The assertion is universally quantified over the four complex parameters and the six algebraic parameters satisfying the hypotheses. The algebraicity assumptions concern the exponential values, not the exponents themselves. No linear independence is required between the two pairs collectively beyond the two separate pairwise hypotheses.

The desired conclusion is exact vanishing in $\mathbb C$, not merely that the exponents are multiples of $2\pi i$ or that their exponentials equal one. That distinction is part of the strength of the formulation. A counterexample would be admissible data satisfying all five algebraicity conditions with at least one displayed exponent nonzero.
\subsection{Short English statement}
Under the two independence hypotheses, do five specified algebraic exponential values force all their exponents to vanish?
\subsection{Sources}
\begin{itemize}
\item[S1] ULAM.AI and the UnsolvedMath Contributors, supplied problems.json, record 8700023 (AMR-086-0023), AMR Open Problem Lists. Catalogue identity and source transcription; CC BY 4.0.
\url{https://huggingface.co/datasets/ulamai/UnsolvedMath}.
\item[S2] Waldschmidt - Open Diophantine Problems (2004), Conjecture 3.5. Original source indexed by AMR-086-0023.
\url{https://arxiv.org/abs/math/0312440}.
\end{itemize}
\end{document}
